%% kapa: exam=1 round=37 year=2026 period=2 form=A subject=appraisal_law subject_ko=감정평가관계법규 seq=12 number=12
%% source: Q-Net 감정평가사 기출문제 (cst00309) artlSeq=5258924 file=raw/1008/5258924/2026년 제37회 감정평가사 1차 2교시.pdf
%% conversion: kordoc 4.24.51 -> tools/gen1.py (auto); review=proofread-sonnet; flags=none
\begin{question}{12}
국토의 계획 및 이용에 관한 법령상 도시ㆍ군계획시설사업의 시행자(이하 ‘시행자’)에 관한 설명으로 옳지 않은 것은?
\begin{choices}
  \item ｢한국관광공사법｣에 따른 한국관광공사는 국토교통부장관, 시ㆍ도지사, 시장 또는 군수로부터 시행자로 지정을 받아 도시ㆍ군계획시설사업을 시행할 수 있다.
  \item 국토교통부장관은 국가계획과 관련된 경우에는 직접 도시ㆍ군계획시설사업을 시행할 수 있다.
  \item ｢대한석탄공사법｣에 따른 대한석탄공사가 시행자로 지정을 받으려면 도시ㆍ군계획시설사업의 대상인 토지면적의 3분의 2 이상에 해당하는 토지를 소유하여야 한다.
  \item 시행자는 도시ㆍ군계획시설사업을 효율적으로 추진하기 위하여 필요하다고 인정되면 사업시행대상지역을 둘 이상으로 분할하여 도시ㆍ군계획시설사업을 시행할 수 있다.
  \item 행정청이 아닌 시행자의 처분에 대하여 행정심판을 제기하는 경우 그 시행자를 지정한 자에게 행정심판을 제기하여야 한다.
\end{choices}
\end{question}
\answer{3}
